\section{Evaluation}\label{sec:eval}

We ask seven questions. \textbf{RQ1:} Does compiling against the data find mistakes before any data
is processed? \textbf{RQ2:} What does data-aware compilation cost and gain? \textbf{RQ3:} Can a
compiled program be reused on edited batches? \textbf{RQ4:} On real data that changes from month to
month, compared with data-validation tools? \textbf{RQ5:} For real pandas scripts? \textbf{RQ6:} Do
certificates compose across stages? \textbf{RQ7:} What do certification and reuse cost? We also report
how we tested soundness and the bugs the evaluation found in our own implementation.

\paragraph{Setup.} Intel Core i5-12450H (8 cores), 16\,GB RAM, Windows~11; MSVC 14.44 (Release);
Python 3.14, pandas 3.0, NumPy 2.5, OpenCV 5.0, pyarrow 25. pandera 0.34 and Great Expectations
0.18.22, whose onboarding data assistant builds a suite from a sample (we found no such assistant in
the 1.x line), run in a Python 3.11 environment, TFDV 1.21 in WSL (Ubuntu, Python 3.12), since it has
no Windows build. Times include starting Python.

\paragraph{Scoring reuse.} RQ3--RQ6 score every way of running a new batch $T$ against the ground
truth: a fresh compilation of the pipeline for $T$, run. A policy either runs a program on $T$ or
refuses it. Its outcome is \emph{correct} if it wrote exactly the files of the fresh run or stopped
where the fresh run stops (a refusal, or a failing check, when the fresh compilation fails or its run
stops); \emph{needless} if it refused $T$ although the fresh compilation runs, which costs a
recompilation; and \emph{wrong} if it finished with other files than the fresh run, finished where the
fresh run fails, or stopped where it does not. The policies are: the \emph{fingerprint} of the
compiled file; the compiled file's own facts (types, no missing values, and the range of every column
read, a schema as TFDV would record it with ranges); pandera's \texttt{infer\_schema} (dtypes,
nullability, numeric ranges); Great Expectations' onboarding data assistant (an inferred suite of
column statistics); TFDV's \texttt{infer\_schema} (types, presence, string domains); the \emph{full
certificate}; and \emph{selective deoptimization}. The validators are set up on the file the pipeline
was compiled with, once on every column (as they are used) and once on the columns the program reads
(more favourable to them); when a validator admits $T$, the optimized program runs on it as it is.

\subsection{RQ1: Mistakes found before any data is processed}

\paragraph{Injected mistakes.} We injected 24 mistakes into the five example programs, one per
mutant~\cite{jia2011}: misspelled columns, type errors (normalizing text, comparing a number with
text), uses after removal, broken contracts (range, not-null, foreign key, schema, channels, width,
match rate), a broken image, and three well-typed programs with silently wrong output. \PL reports all
24 at compile time and nothing on the unmodified programs, in \RqOneMsMin--\RqOneMsMax\,ms; a
script-style interpreter of the same programs crashes part-way on 20 and finishes 4 with wrong output
(text filled into an integer column, a threshold on a normalized column that removes every row, and
two pipelines that write black images). On 16 million Parquet rows, a misspelled column is reported in
\ScTabNamePl\,s instead of after \ScTabNameScript\,s, and float pixels saved as PNG in \ScImgFloatPl\,s
for 3,000 images, where the script runs \ScImgFloatScript\,s and writes black images.

\paragraph{Real failures.} We chose 17 error messages that Python libraries print for common
data-preparation failures, searched GitHub issues for each and read the first 12 results in order; an
issue qualified if the failure happened while loading, transforming or saving data files and its root
cause was stated, at most three per message: \RbTotal{} issues~\cite{islam2019,humbatova2020}. \PL
reports \RbEarly{} of them (\RbEarlyPct\%) before any data is processed (\RbCompile{} without
annotation, \RbContract{} once the requirement is declared with \texttt{EXPECT}); \RbDesign{} cannot
occur in a \PL program, and \RbNone{} are not covered (input formats such as HEIC or GeoTIFF,
operations such as pivots, dates, per-image label files). \RbPass{} of \RbCases{} reproductions behave
as classified.

\subsection{RQ2: Cost and gain of data-aware compilation}

Analysing every cell of a one-million-row CSV file takes \CsvMillion\,s, \CsvSpeedup$\times$ faster
than pandas' default reader reads it; Parquet compilation stays at \PqSixteen\,s from one to 16
million rows; reading 3,000 image headers is \HdrRatio$\times$ faster than decoding the images. On a
200,000-row, 40-column table with two contracts and a filter, the optimizer proves both contracts and
loads 5 of 40 columns; the optimized program processes the data in \TabProcOpt\,s instead of
\TabProcOz\,s (\TabCut\% less) with byte-identical output. Image pipelines stream one image at a
time, in \ImgPlMem\,MiB of memory instead of \ImgBaseMem\,MiB for loading all 1,000 images.

\subsection{RQ3: Reusing a compiled program on edited batches}\label{sec:rq-edits}

Seven workloads are compiled once: the four table examples, one program run twice on one file, the
Parquet example, images linked to labels, and a sensor pipeline of 20,000 rows per day. Each input
table gets new batches: other days and a warmer season of the sensor process, and, for the examples,
rows resampled with replacement, values moved by noise of 0.1, 0.5 and 1.5 standard deviations, 20\%
missing values in one column, and one value far out of range (\CertBatches{} batches).

\begin{table*}[t]
\caption{Reusing compiled programs on edited batches (RQ3): per policy, batches handled correctly (ok: the fresh files, or a stop where the fresh run stops), refused although a fresh compilation runs (ref.), and finished with other files or where the fresh run fails (wrong).}
\label{tab:rq-edits}
\centering
\footnotesize
\begin{adjustbox}{max width=\linewidth}
\begin{tabular}{@{}lrrrrrrrrr@{}}
\toprule
& \multicolumn{3}{c}{tables} & \multicolumn{3}{c}{images + labels} & \multicolumn{3}{c}{all} \\
\cmidrule(l){2-4}\cmidrule(l){5-7}\cmidrule(l){8-10}
Policy & ok & ref. & wrong & ok & ref. & wrong & ok & ref. & wrong \\
\midrule
Fingerprint & 98 & 126 & 0 & 17 & 3 & 0 & 115 & 129 & 0 \\
File facts & 121 & 78 & 25 & 1 & 3 & 16 & 122 & 81 & 41 \\
pandera, every column & 120 & 79 & 25 & 1 & 3 & 16 & 121 & 82 & 41 \\
pandera, columns read & 123 & 76 & 25 & 1 & 3 & 16 & 124 & 79 & 41 \\
Great Expectations, every column & 121 & 103 & 0 & 17 & 3 & 0 & 138 & 106 & 0 \\
Great Expectations, columns read & 123 & 101 & 0 & 17 & 3 & 0 & 140 & 104 & 0 \\
TFDV, every column & 146 & 28 & 50 & 1 & 3 & 16 & 147 & 31 & 66 \\
TFDV, columns read & 149 & 25 & 50 & 1 & 3 & 16 & 150 & 28 & 66 \\
Full certificate & 224 & 0 & 0 & 20 & 0 & 0 & 244 & 0 & 0 \\
\textbf{Selective deoptimization} & \textbf{224} & \textbf{0} & \textbf{0} & \textbf{20} & \textbf{0} & \textbf{0} & \textbf{244} & \textbf{0} & \textbf{0} \\
\bottomrule
\end{tabular}
\end{adjustbox}
\end{table*}


Both certificates handle every batch correctly (Table~\ref{tab:rq-edits}): the full certificate by
refusing exactly the batches whose fresh compilation fails, selective deoptimization by running them
and stopping at the re-enabled check (\CertDeopted{} batches ran with removals re-enabled). The
fingerprint refuses every edited batch (\CertFpNeedless{} of them needlessly), and the compiled file's
own facts refuse \CertFactsNeedless{} that a fresh compilation runs. Gating the
optimized program with the file's own facts is also unsound (\CertFactsUnsound{} wrong runs): on the
table set and on the linked images, resampled and drifted rows repeat keys or lose matches, so a
contract the compiler had proven by reading the values (\texttt{UNIQUE}, \texttt{REFERENCES},
\texttt{MATCH\_RATE}) no longer holds and the fresh compilation reports it, while the file's facts,
which are facts of single columns, admit the batch and the optimized program finishes. \PL runs such
checks again whenever an input changed. The real validators fare no better: pandera's inferred schema
is unsound on the same batches (\CertPanderaAllUnsound{} wrong), TFDV's, which has no ranges, on these
and on batches whose values drift past what a removed check or filter relied on (\CertTfdvAllUnsound{}
wrong), and Great Expectations' statistical suite, which
is sound here, refuses \CertGxAllNeedless{} batches that a fresh compilation runs.

\paragraph{Image folders.} A pipeline that makes images gray and $64\times64$, compiled on a folder
of gray $64\times64$ images, has its \texttt{GRAYSCALE} and \texttt{RESIZE} removed speculatively. We
gave it \ImgBatches{} new folders: gray $64\times64$ images, some color images, other sizes, both, a
file that is not an image, and 16-bit images (five each), and compared the written images byte by
byte with a fresh compilation. Selective deoptimization handles \ImgBDeoptCorrect{} correctly,
re-running \texttt{GRAYSCALE} or \texttt{RESIZE} on \ImgBDeopted{} folders; the full certificate
refuses those folders, and the optimized program run on them without a check writes color or
wrongly sized images (\ImgBTrustedUnsound{} folders). The \ImgBDeoptNeedless{} needless refusals are
the 16-bit folders: bit depths are fixed in every certificate because they decide the pixel type the
generated code assumes, even where, as here, no instruction depends on it; certifying them would need
the code generator's assumptions as obligations too.

\subsection{RQ4: Real monthly data, against data-validation tools}\label{sec:rq-real}

Three public datasets come in monthly batches: hourly traffic and weather on a Minnesota interstate
(UCI, 2012--2018, 63 months of CSV), daily bike rentals in Washington (UCI, 2011--2012, 24 months),
and New York green-taxi trips (NYC TLC, 2022--2024, the 36 Parquet files as published). For each we
wrote one cleaning pipeline as a data engineer would: range contracts on the measured quantities,
\texttt{NOT\_NULL} on keys, fills, a filter that removes impossible readings, a normalization and a
selection. Each pipeline is deployed twice: compiled on the first month and run on every later one,
and compiled on the month before each month (\RealRuns{} runs; twice the month before has a temperature of
0\,K, fails to compile itself, and leaves nothing to deploy). Batches are cut from the raw
files, so no value is reformatted.

\begin{table*}[t]
\caption{Reusing compiled pipelines on real monthly batches (RQ4), both deployments (compiled on the first month, and on the month before). ok: handled as a fresh compilation of the month would; ref.: refused although a fresh compilation runs; wrong: finished with other files than the fresh run, or where it fails. Validators gate the optimized program (Section~\ref{sec:rq-real}).}
\label{tab:rq-real}
\centering
\footnotesize
\begin{adjustbox}{max width=\linewidth}
\begin{tabular}{@{}lrrrrrrrrrrrrrrrrrr@{}}
\toprule
& \multicolumn{3}{c}{Traffic (CSV)} & \multicolumn{3}{c}{Bikes (CSV)} & \multicolumn{3}{c}{Taxi (Parquet)} & \multicolumn{3}{c}{Citi Bike (CSV)} & \multicolumn{3}{c}{Air quality (CSV)} & \multicolumn{3}{c}{all 426} \\
\cmidrule(l){2-4}\cmidrule(l){5-7}\cmidrule(l){8-10}\cmidrule(l){11-13}\cmidrule(l){14-16}\cmidrule(l){17-19}
Policy & ok & ref. & wrong & ok & ref. & wrong & ok & ref. & wrong & ok & ref. & wrong & ok & ref. & wrong & ok & ref. & wrong \\
\midrule
Fingerprint & 3 & 119 & 0 & 0 & 46 & 0 & 14 & 56 & 0 & 0 & 70 & 0 & 0 & 118 & 0 & 17 & 409 & 0 \\
File facts & 5 & 117 & 0 & 0 & 46 & 0 & 14 & 56 & 0 & 0 & 70 & 0 & 0 & 118 & 0 & 19 & 407 & 0 \\
pandera, every column & 5 & 117 & 0 & 0 & 46 & 0 & 14 & 56 & 0 & 0 & 70 & 0 & 0 & 118 & 0 & 19 & 407 & 0 \\
pandera, columns read & 5 & 117 & 0 & 0 & 46 & 0 & 14 & 56 & 0 & 0 & 70 & 0 & 0 & 118 & 0 & 19 & 407 & 0 \\
Great Expectations, every column & 3 & 119 & 0 & 0 & 46 & 0 & 14 & 56 & 0 & 0 & 70 & 0 & 0 & 118 & 0 & 17 & 409 & 0 \\
Great Expectations, columns read & 3 & 119 & 0 & 0 & 46 & 0 & 14 & 56 & 0 & 0 & 70 & 0 & 0 & 118 & 0 & 17 & 409 & 0 \\
TFDV, every column & 5 & 117 & 0 & 0 & 46 & 0 & 63 & 0 & 7 & 13 & 57 & 0 & 49 & 69 & 0 & 130 & 289 & 7 \\
TFDV, columns read & 92 & 27 & 3 & 0 & 46 & 0 & 63 & 0 & 7 & 13 & 57 & 0 & 49 & 69 & 0 & 217 & 199 & 10 \\
Full certificate & 120 & 2 & 0 & 44 & 2 & 0 & 68 & 2 & 0 & 50 & 20 & 0 & 115 & 3 & 0 & 397 & 29 & 0 \\
\textbf{Selective deoptimization} & \textbf{122} & \textbf{0} & \textbf{0} & \textbf{46} & \textbf{0} & \textbf{0} & \textbf{70} & \textbf{0} & \textbf{0} & \textbf{70} & \textbf{0} & \textbf{0} & \textbf{118} & \textbf{0} & \textbf{0} & \textbf{426} & \textbf{0} & \textbf{0} \\
\bottomrule
\end{tabular}
\end{adjustbox}
\end{table*}

\begin{figure}[t]
\centering
\begin{adjustbox}{max width=\linewidth}
\begin{tikzpicture}
\begin{axis}[
    xbar stacked, bar width=7pt, width=\columnwidth, height=6.2cm,
    xmin=0, xmax=\RealRuns, enlarge y limits=0.06,
    xlabel={real monthly batches}, xlabel style={font=\scriptsize},
    symbolic y coords={TFDV read,TFDV all,GX read,GX all,pandera read,pandera all,File facts,Fingerprint,Full cert.,Deopt.},
    ytick=data, yticklabel style={font=\scriptsize}, xticklabel style={font=\scriptsize},
    legend style={font=\scriptsize, at={(0.5,1.02)}, anchor=south, legend columns=3, draw=none},
    reverse legend=false]
\addplot[fill=black!70, draw=black!70] coordinates {
    (\RealTfdvReadCorrect,TFDV read) (\RealTfdvAllCorrect,TFDV all) (\RealGxReadCorrect,GX read)
    (\RealGxAllCorrect,GX all) (\RealPanderaReadCorrect,pandera read) (\RealPanderaAllCorrect,pandera all)
    (\RealFactsCorrect,File facts) (\RealFpCorrect,Fingerprint) (\RealFullCorrect,Full cert.)
    (\RealDeoptCorrect,Deopt.)};
\addplot[fill=black!20, draw=black!40] coordinates {
    (\RealTfdvReadNeedless,TFDV read) (\RealTfdvAllNeedless,TFDV all) (\RealGxReadNeedless,GX read)
    (\RealGxAllNeedless,GX all) (\RealPanderaReadNeedless,pandera read) (\RealPanderaAllNeedless,pandera all)
    (\RealFactsNeedless,File facts) (\RealFpNeedless,Fingerprint) (\RealFullNeedless,Full cert.)
    (\RealDeoptNeedless,Deopt.)};
\addplot[pattern=north east lines, pattern color=black, draw=black] coordinates {
    (\RealTfdvReadUnsound,TFDV read) (\RealTfdvAllUnsound,TFDV all) (\RealGxReadUnsound,GX read)
    (\RealGxAllUnsound,GX all) (\RealPanderaReadUnsound,pandera read) (\RealPanderaAllUnsound,pandera all)
    (\RealFactsUnsound,File facts) (\RealFpUnsound,Fingerprint) (\RealFullUnsound,Full cert.)
    (\RealDeoptUnsound,Deopt.)};
\legend{correct, refused needlessly, wrong}
\end{axis}
\end{tikzpicture}
\end{adjustbox}
\caption{RQ4: \RealRuns{} real monthly batches, by admission policy. \emph{Correct}: ran and wrote
the fresh compilation's files, or stopped where it stops; \emph{refused needlessly}: a recompilation
was needed although the fresh compilation runs; \emph{wrong}: other files, or a run where the fresh
compilation fails. Validators gate the optimized program; \emph{all}/\emph{read}: set up on every
column, or on the columns the program reads.}
\label{fig:policies}
\end{figure}

\begin{figure*}[t]
\centering
\definecolor{deopt}{RGB}{217,95,2}
\begin{adjustbox}{max width=\linewidth}
\begin{tikzpicture}[font=\scriptsize]
\node[anchor=west, inner sep=0pt, font=\scriptsize\bfseries] at (0,0.000) {Traffic};
\node[anchor=west, inner sep=0pt, text=black!60] at (1.980,0.000) {2013};
\node[anchor=west, inner sep=0pt, text=black!60] at (4.560,0.000) {2014};
\node[anchor=west, inner sep=0pt, text=black!60] at (7.140,0.000) {2015};
\node[anchor=west, inner sep=0pt, text=black!60] at (9.720,0.000) {2016};
\node[anchor=west, inner sep=0pt, text=black!60] at (12.300,0.000) {2017};
\node[anchor=west, inner sep=0pt, text=black!60] at (14.880,0.000) {2018};
\node[anchor=east] at (1.470,-0.273) {TFDV read};
\path[fill=black!65, draw=black!65] (1.550,-0.380) rectangle ++(0.180,0.215);
\path[fill=black!65, draw=black!65] (1.765,-0.380) rectangle ++(0.180,0.215);
\path[fill=black!65, draw=black!65] (1.980,-0.380) rectangle ++(0.180,0.215);
\path[fill=black!65, draw=black!65] (2.195,-0.380) rectangle ++(0.180,0.215);
\path[fill=black!65, draw=black!65] (2.410,-0.380) rectangle ++(0.180,0.215);
\path[fill=black!65, draw=black!65] (2.625,-0.380) rectangle ++(0.180,0.215);
\path[fill=black!15, draw=black!15] (2.840,-0.380) rectangle ++(0.180,0.215);
\path[fill=black!15, draw=black!15] (3.055,-0.380) rectangle ++(0.180,0.215);
\path[fill=black!65, draw=black!65] (3.270,-0.380) rectangle ++(0.180,0.215);
\path[fill=black!15, draw=black!15] (3.485,-0.380) rectangle ++(0.180,0.215);
\path[fill=black!65, draw=black!65] (3.700,-0.380) rectangle ++(0.180,0.215);
\path[fill=black!65, draw=black!65] (3.915,-0.380) rectangle ++(0.180,0.215);
\path[fill=black!65, draw=black!65] (4.130,-0.380) rectangle ++(0.180,0.215);
\path[fill=black!65, draw=black!65] (4.345,-0.380) rectangle ++(0.180,0.215);
\path[pattern=north east lines, pattern color=black, draw=black] (4.560,-0.380) rectangle ++(0.180,0.215);
\path[pattern=north east lines, pattern color=black, draw=black] (4.775,-0.380) rectangle ++(0.180,0.215);
\path[fill=black!65, draw=black!65] (4.990,-0.380) rectangle ++(0.180,0.215);
\path[fill=black!65, draw=black!65] (5.205,-0.380) rectangle ++(0.180,0.215);
\path[fill=black!65, draw=black!65] (5.420,-0.380) rectangle ++(0.180,0.215);
\path[fill=black!65, draw=black!65] (5.635,-0.380) rectangle ++(0.180,0.215);
\path[fill=black!65, draw=black!65] (5.850,-0.380) rectangle ++(0.180,0.215);
\path[fill=black!65, draw=black!65] (6.065,-0.380) rectangle ++(0.180,0.215);
\path[fill=black!65, draw=black!65] (8.215,-0.380) rectangle ++(0.180,0.215);
\path[fill=black!15, draw=black!15] (8.430,-0.380) rectangle ++(0.180,0.215);
\path[fill=black!65, draw=black!65] (8.645,-0.380) rectangle ++(0.180,0.215);
\path[fill=black!65, draw=black!65] (8.860,-0.380) rectangle ++(0.180,0.215);
\path[fill=black!65, draw=black!65] (9.075,-0.380) rectangle ++(0.180,0.215);
\path[fill=black!65, draw=black!65] (9.290,-0.380) rectangle ++(0.180,0.215);
\path[fill=black!65, draw=black!65] (9.505,-0.380) rectangle ++(0.180,0.215);
\path[fill=black!65, draw=black!65] (9.720,-0.380) rectangle ++(0.180,0.215);
\path[fill=black!65, draw=black!65] (9.935,-0.380) rectangle ++(0.180,0.215);
\path[fill=black!65, draw=black!65] (10.150,-0.380) rectangle ++(0.180,0.215);
\path[fill=black!65, draw=black!65] (10.365,-0.380) rectangle ++(0.180,0.215);
\path[fill=black!15, draw=black!15] (10.580,-0.380) rectangle ++(0.180,0.215);
\path[fill=black!65, draw=black!65] (10.795,-0.380) rectangle ++(0.180,0.215);
\path[fill=black!15, draw=black!15] (11.010,-0.380) rectangle ++(0.180,0.215);
\path[fill=black!65, draw=black!65] (11.225,-0.380) rectangle ++(0.180,0.215);
\path[fill=black!65, draw=black!65] (11.440,-0.380) rectangle ++(0.180,0.215);
\path[fill=black!65, draw=black!65] (11.655,-0.380) rectangle ++(0.180,0.215);
\path[fill=black!65, draw=black!65] (11.870,-0.380) rectangle ++(0.180,0.215);
\path[fill=black!65, draw=black!65] (12.085,-0.380) rectangle ++(0.180,0.215);
\path[fill=black!65, draw=black!65] (12.300,-0.380) rectangle ++(0.180,0.215);
\path[fill=black!65, draw=black!65] (12.515,-0.380) rectangle ++(0.180,0.215);
\path[fill=black!65, draw=black!65] (12.730,-0.380) rectangle ++(0.180,0.215);
\path[fill=black!65, draw=black!65] (12.945,-0.380) rectangle ++(0.180,0.215);
\path[fill=black!65, draw=black!65] (13.160,-0.380) rectangle ++(0.180,0.215);
\path[fill=black!65, draw=black!65] (13.375,-0.380) rectangle ++(0.180,0.215);
\path[fill=black!65, draw=black!65] (13.590,-0.380) rectangle ++(0.180,0.215);
\path[fill=black!65, draw=black!65] (13.805,-0.380) rectangle ++(0.180,0.215);
\path[fill=black!15, draw=black!15] (14.020,-0.380) rectangle ++(0.180,0.215);
\path[fill=black!65, draw=black!65] (14.235,-0.380) rectangle ++(0.180,0.215);
\path[fill=black!65, draw=black!65] (14.450,-0.380) rectangle ++(0.180,0.215);
\path[fill=black!65, draw=black!65] (14.665,-0.380) rectangle ++(0.180,0.215);
\path[fill=black!65, draw=black!65] (14.880,-0.380) rectangle ++(0.180,0.215);
\path[fill=black!65, draw=black!65] (15.095,-0.380) rectangle ++(0.180,0.215);
\path[fill=black!65, draw=black!65] (15.310,-0.380) rectangle ++(0.180,0.215);
\path[fill=black!65, draw=black!65] (15.525,-0.380) rectangle ++(0.180,0.215);
\path[fill=black!65, draw=black!65] (15.740,-0.380) rectangle ++(0.180,0.215);
\path[fill=black!65, draw=black!65] (15.955,-0.380) rectangle ++(0.180,0.215);
\path[fill=black!65, draw=black!65] (16.170,-0.380) rectangle ++(0.180,0.215);
\path[fill=black!15, draw=black!15] (16.385,-0.380) rectangle ++(0.180,0.215);
\path[fill=black!65, draw=black!65] (16.600,-0.380) rectangle ++(0.180,0.215);
\node[anchor=east] at (1.470,-0.522) {Full cert.};
\path[fill=black!65, draw=black!65] (1.550,-0.630) rectangle ++(0.180,0.215);
\path[fill=black!65, draw=black!65] (1.765,-0.630) rectangle ++(0.180,0.215);
\path[fill=black!65, draw=black!65] (1.980,-0.630) rectangle ++(0.180,0.215);
\path[fill=black!65, draw=black!65] (2.195,-0.630) rectangle ++(0.180,0.215);
\path[fill=black!65, draw=black!65] (2.410,-0.630) rectangle ++(0.180,0.215);
\path[fill=black!65, draw=black!65] (2.625,-0.630) rectangle ++(0.180,0.215);
\path[fill=black!65, draw=black!65] (2.840,-0.630) rectangle ++(0.180,0.215);
\path[fill=black!65, draw=black!65] (3.055,-0.630) rectangle ++(0.180,0.215);
\path[fill=black!65, draw=black!65] (3.270,-0.630) rectangle ++(0.180,0.215);
\path[fill=black!65, draw=black!65] (3.485,-0.630) rectangle ++(0.180,0.215);
\path[fill=black!65, draw=black!65] (3.700,-0.630) rectangle ++(0.180,0.215);
\path[fill=black!65, draw=black!65] (3.915,-0.630) rectangle ++(0.180,0.215);
\path[fill=black!65, draw=black!65] (4.130,-0.630) rectangle ++(0.180,0.215);
\path[fill=black!65, draw=black!65] (4.345,-0.630) rectangle ++(0.180,0.215);
\path[fill=white, draw=black, line width=0.5pt] (4.560,-0.630) rectangle ++(0.180,0.215);
\path[fill=white, draw=black, line width=0.5pt] (4.775,-0.630) rectangle ++(0.180,0.215);
\path[fill=black!65, draw=black!65] (4.990,-0.630) rectangle ++(0.180,0.215);
\path[fill=black!65, draw=black!65] (5.205,-0.630) rectangle ++(0.180,0.215);
\path[fill=black!65, draw=black!65] (5.420,-0.630) rectangle ++(0.180,0.215);
\path[fill=black!65, draw=black!65] (5.635,-0.630) rectangle ++(0.180,0.215);
\path[fill=black!65, draw=black!65] (5.850,-0.630) rectangle ++(0.180,0.215);
\path[fill=black!65, draw=black!65] (6.065,-0.630) rectangle ++(0.180,0.215);
\path[fill=black!65, draw=black!65] (8.215,-0.630) rectangle ++(0.180,0.215);
\path[fill=black!65, draw=black!65] (8.430,-0.630) rectangle ++(0.180,0.215);
\path[fill=black!65, draw=black!65] (8.645,-0.630) rectangle ++(0.180,0.215);
\path[fill=black!65, draw=black!65] (8.860,-0.630) rectangle ++(0.180,0.215);
\path[fill=black!65, draw=black!65] (9.075,-0.630) rectangle ++(0.180,0.215);
\path[fill=black!65, draw=black!65] (9.290,-0.630) rectangle ++(0.180,0.215);
\path[fill=black!65, draw=black!65] (9.505,-0.630) rectangle ++(0.180,0.215);
\path[fill=black!65, draw=black!65] (9.720,-0.630) rectangle ++(0.180,0.215);
\path[fill=black!65, draw=black!65] (9.935,-0.630) rectangle ++(0.180,0.215);
\path[fill=black!65, draw=black!65] (10.150,-0.630) rectangle ++(0.180,0.215);
\path[fill=black!65, draw=black!65] (10.365,-0.630) rectangle ++(0.180,0.215);
\path[fill=black!65, draw=black!65] (10.580,-0.630) rectangle ++(0.180,0.215);
\path[fill=black!65, draw=black!65] (10.795,-0.630) rectangle ++(0.180,0.215);
\path[fill=black!15, draw=black!15] (11.010,-0.630) rectangle ++(0.180,0.215);
\path[fill=black!65, draw=black!65] (11.225,-0.630) rectangle ++(0.180,0.215);
\path[fill=black!65, draw=black!65] (11.440,-0.630) rectangle ++(0.180,0.215);
\path[fill=black!65, draw=black!65] (11.655,-0.630) rectangle ++(0.180,0.215);
\path[fill=black!65, draw=black!65] (11.870,-0.630) rectangle ++(0.180,0.215);
\path[fill=black!65, draw=black!65] (12.085,-0.630) rectangle ++(0.180,0.215);
\path[fill=black!65, draw=black!65] (12.300,-0.630) rectangle ++(0.180,0.215);
\path[fill=black!65, draw=black!65] (12.515,-0.630) rectangle ++(0.180,0.215);
\path[fill=black!65, draw=black!65] (12.730,-0.630) rectangle ++(0.180,0.215);
\path[fill=black!65, draw=black!65] (12.945,-0.630) rectangle ++(0.180,0.215);
\path[fill=black!65, draw=black!65] (13.160,-0.630) rectangle ++(0.180,0.215);
\path[fill=black!65, draw=black!65] (13.375,-0.630) rectangle ++(0.180,0.215);
\path[fill=black!65, draw=black!65] (13.590,-0.630) rectangle ++(0.180,0.215);
\path[fill=black!65, draw=black!65] (13.805,-0.630) rectangle ++(0.180,0.215);
\path[fill=black!65, draw=black!65] (14.020,-0.630) rectangle ++(0.180,0.215);
\path[fill=black!65, draw=black!65] (14.235,-0.630) rectangle ++(0.180,0.215);
\path[fill=black!65, draw=black!65] (14.450,-0.630) rectangle ++(0.180,0.215);
\path[fill=black!65, draw=black!65] (14.665,-0.630) rectangle ++(0.180,0.215);
\path[fill=black!65, draw=black!65] (14.880,-0.630) rectangle ++(0.180,0.215);
\path[fill=black!65, draw=black!65] (15.095,-0.630) rectangle ++(0.180,0.215);
\path[fill=black!65, draw=black!65] (15.310,-0.630) rectangle ++(0.180,0.215);
\path[fill=black!65, draw=black!65] (15.525,-0.630) rectangle ++(0.180,0.215);
\path[fill=black!65, draw=black!65] (15.740,-0.630) rectangle ++(0.180,0.215);
\path[fill=black!65, draw=black!65] (15.955,-0.630) rectangle ++(0.180,0.215);
\path[fill=black!65, draw=black!65] (16.170,-0.630) rectangle ++(0.180,0.215);
\path[fill=black!65, draw=black!65] (16.385,-0.630) rectangle ++(0.180,0.215);
\path[fill=black!65, draw=black!65] (16.600,-0.630) rectangle ++(0.180,0.215);
\node[anchor=east] at (1.470,-0.772) {Deopt.};
\path[fill=black!65, draw=black!65] (1.550,-0.880) rectangle ++(0.180,0.215);
\path[fill=black!65, draw=black!65] (1.765,-0.880) rectangle ++(0.180,0.215);
\path[fill=black!65, draw=black!65] (1.980,-0.880) rectangle ++(0.180,0.215);
\path[fill=black!65, draw=black!65] (2.195,-0.880) rectangle ++(0.180,0.215);
\path[fill=black!65, draw=black!65] (2.410,-0.880) rectangle ++(0.180,0.215);
\path[fill=black!65, draw=black!65] (2.625,-0.880) rectangle ++(0.180,0.215);
\path[fill=black!65, draw=black!65] (2.840,-0.880) rectangle ++(0.180,0.215);
\path[fill=black!65, draw=black!65] (3.055,-0.880) rectangle ++(0.180,0.215);
\path[fill=black!65, draw=black!65] (3.270,-0.880) rectangle ++(0.180,0.215);
\path[fill=black!65, draw=black!65] (3.485,-0.880) rectangle ++(0.180,0.215);
\path[fill=black!65, draw=black!65] (3.700,-0.880) rectangle ++(0.180,0.215);
\path[fill=black!65, draw=black!65] (3.915,-0.880) rectangle ++(0.180,0.215);
\path[fill=black!65, draw=black!65] (4.130,-0.880) rectangle ++(0.180,0.215);
\path[fill=black!65, draw=black!65] (4.345,-0.880) rectangle ++(0.180,0.215);
\path[fill=white, draw=black, line width=0.5pt] (4.560,-0.880) rectangle ++(0.180,0.215);
\path[fill=white, draw=black, line width=0.5pt] (4.775,-0.880) rectangle ++(0.180,0.215);
\path[fill=black!65, draw=black!65] (4.990,-0.880) rectangle ++(0.180,0.215);
\path[fill=black!65, draw=black!65] (5.205,-0.880) rectangle ++(0.180,0.215);
\path[fill=black!65, draw=black!65] (5.420,-0.880) rectangle ++(0.180,0.215);
\path[fill=black!65, draw=black!65] (5.635,-0.880) rectangle ++(0.180,0.215);
\path[fill=black!65, draw=black!65] (5.850,-0.880) rectangle ++(0.180,0.215);
\path[fill=black!65, draw=black!65] (6.065,-0.880) rectangle ++(0.180,0.215);
\path[fill=black!65, draw=black!65] (8.215,-0.880) rectangle ++(0.180,0.215);
\path[fill=black!65, draw=black!65] (8.430,-0.880) rectangle ++(0.180,0.215);
\path[fill=black!65, draw=black!65] (8.645,-0.880) rectangle ++(0.180,0.215);
\path[fill=black!65, draw=black!65] (8.860,-0.880) rectangle ++(0.180,0.215);
\path[fill=black!65, draw=black!65] (9.075,-0.880) rectangle ++(0.180,0.215);
\path[fill=black!65, draw=black!65] (9.290,-0.880) rectangle ++(0.180,0.215);
\path[fill=black!65, draw=black!65] (9.505,-0.880) rectangle ++(0.180,0.215);
\path[fill=black!65, draw=black!65] (9.720,-0.880) rectangle ++(0.180,0.215);
\path[fill=black!65, draw=black!65] (9.935,-0.880) rectangle ++(0.180,0.215);
\path[fill=black!65, draw=black!65] (10.150,-0.880) rectangle ++(0.180,0.215);
\path[fill=black!65, draw=black!65] (10.365,-0.880) rectangle ++(0.180,0.215);
\path[fill=black!65, draw=black!65] (10.580,-0.880) rectangle ++(0.180,0.215);
\path[fill=black!65, draw=black!65] (10.795,-0.880) rectangle ++(0.180,0.215);
\path[fill=deopt, draw=deopt] (11.010,-0.880) rectangle ++(0.180,0.215);
\path[fill=black!65, draw=black!65] (11.225,-0.880) rectangle ++(0.180,0.215);
\path[fill=black!65, draw=black!65] (11.440,-0.880) rectangle ++(0.180,0.215);
\path[fill=black!65, draw=black!65] (11.655,-0.880) rectangle ++(0.180,0.215);
\path[fill=black!65, draw=black!65] (11.870,-0.880) rectangle ++(0.180,0.215);
\path[fill=black!65, draw=black!65] (12.085,-0.880) rectangle ++(0.180,0.215);
\path[fill=black!65, draw=black!65] (12.300,-0.880) rectangle ++(0.180,0.215);
\path[fill=black!65, draw=black!65] (12.515,-0.880) rectangle ++(0.180,0.215);
\path[fill=black!65, draw=black!65] (12.730,-0.880) rectangle ++(0.180,0.215);
\path[fill=black!65, draw=black!65] (12.945,-0.880) rectangle ++(0.180,0.215);
\path[fill=black!65, draw=black!65] (13.160,-0.880) rectangle ++(0.180,0.215);
\path[fill=black!65, draw=black!65] (13.375,-0.880) rectangle ++(0.180,0.215);
\path[fill=black!65, draw=black!65] (13.590,-0.880) rectangle ++(0.180,0.215);
\path[fill=black!65, draw=black!65] (13.805,-0.880) rectangle ++(0.180,0.215);
\path[fill=black!65, draw=black!65] (14.020,-0.880) rectangle ++(0.180,0.215);
\path[fill=black!65, draw=black!65] (14.235,-0.880) rectangle ++(0.180,0.215);
\path[fill=black!65, draw=black!65] (14.450,-0.880) rectangle ++(0.180,0.215);
\path[fill=black!65, draw=black!65] (14.665,-0.880) rectangle ++(0.180,0.215);
\path[fill=black!65, draw=black!65] (14.880,-0.880) rectangle ++(0.180,0.215);
\path[fill=black!65, draw=black!65] (15.095,-0.880) rectangle ++(0.180,0.215);
\path[fill=black!65, draw=black!65] (15.310,-0.880) rectangle ++(0.180,0.215);
\path[fill=black!65, draw=black!65] (15.525,-0.880) rectangle ++(0.180,0.215);
\path[fill=black!65, draw=black!65] (15.740,-0.880) rectangle ++(0.180,0.215);
\path[fill=black!65, draw=black!65] (15.955,-0.880) rectangle ++(0.180,0.215);
\path[fill=black!65, draw=black!65] (16.170,-0.880) rectangle ++(0.180,0.215);
\path[fill=black!65, draw=black!65] (16.385,-0.880) rectangle ++(0.180,0.215);
\path[fill=black!65, draw=black!65] (16.600,-0.880) rectangle ++(0.180,0.215);
\node[anchor=west, inner sep=0pt, font=\scriptsize\bfseries] at (0,-1.300) {Bikes};
\node[anchor=west, inner sep=0pt, text=black!60] at (1.550,-1.300) {2011};
\node[anchor=west, inner sep=0pt, text=black!60] at (3.915,-1.300) {2012};
\node[anchor=east] at (1.470,-1.573) {TFDV read};
\path[fill=black!15, draw=black!15] (1.550,-1.680) rectangle ++(0.180,0.215);
\path[fill=black!15, draw=black!15] (1.765,-1.680) rectangle ++(0.180,0.215);
\path[fill=black!15, draw=black!15] (1.980,-1.680) rectangle ++(0.180,0.215);
\path[fill=black!15, draw=black!15] (2.195,-1.680) rectangle ++(0.180,0.215);
\path[fill=black!15, draw=black!15] (2.410,-1.680) rectangle ++(0.180,0.215);
\path[fill=black!15, draw=black!15] (2.625,-1.680) rectangle ++(0.180,0.215);
\path[fill=black!15, draw=black!15] (2.840,-1.680) rectangle ++(0.180,0.215);
\path[fill=black!15, draw=black!15] (3.055,-1.680) rectangle ++(0.180,0.215);
\path[fill=black!15, draw=black!15] (3.270,-1.680) rectangle ++(0.180,0.215);
\path[fill=black!15, draw=black!15] (3.485,-1.680) rectangle ++(0.180,0.215);
\path[fill=black!15, draw=black!15] (3.700,-1.680) rectangle ++(0.180,0.215);
\path[fill=black!15, draw=black!15] (3.915,-1.680) rectangle ++(0.180,0.215);
\path[fill=black!15, draw=black!15] (4.130,-1.680) rectangle ++(0.180,0.215);
\path[fill=black!15, draw=black!15] (4.345,-1.680) rectangle ++(0.180,0.215);
\path[fill=black!15, draw=black!15] (4.560,-1.680) rectangle ++(0.180,0.215);
\path[fill=black!15, draw=black!15] (4.775,-1.680) rectangle ++(0.180,0.215);
\path[fill=black!15, draw=black!15] (4.990,-1.680) rectangle ++(0.180,0.215);
\path[fill=black!15, draw=black!15] (5.205,-1.680) rectangle ++(0.180,0.215);
\path[fill=black!15, draw=black!15] (5.420,-1.680) rectangle ++(0.180,0.215);
\path[fill=black!15, draw=black!15] (5.635,-1.680) rectangle ++(0.180,0.215);
\path[fill=black!15, draw=black!15] (5.850,-1.680) rectangle ++(0.180,0.215);
\path[fill=black!15, draw=black!15] (6.065,-1.680) rectangle ++(0.180,0.215);
\path[fill=black!15, draw=black!15] (6.280,-1.680) rectangle ++(0.180,0.215);
\node[anchor=east] at (1.470,-1.823) {Full cert.};
\path[fill=black!65, draw=black!65] (1.550,-1.930) rectangle ++(0.180,0.215);
\path[fill=black!15, draw=black!15] (1.765,-1.930) rectangle ++(0.180,0.215);
\path[fill=black!65, draw=black!65] (1.980,-1.930) rectangle ++(0.180,0.215);
\path[fill=black!65, draw=black!65] (2.195,-1.930) rectangle ++(0.180,0.215);
\path[fill=black!65, draw=black!65] (2.410,-1.930) rectangle ++(0.180,0.215);
\path[fill=black!65, draw=black!65] (2.625,-1.930) rectangle ++(0.180,0.215);
\path[fill=black!65, draw=black!65] (2.840,-1.930) rectangle ++(0.180,0.215);
\path[fill=black!65, draw=black!65] (3.055,-1.930) rectangle ++(0.180,0.215);
\path[fill=black!65, draw=black!65] (3.270,-1.930) rectangle ++(0.180,0.215);
\path[fill=black!65, draw=black!65] (3.485,-1.930) rectangle ++(0.180,0.215);
\path[fill=black!65, draw=black!65] (3.700,-1.930) rectangle ++(0.180,0.215);
\path[fill=black!65, draw=black!65] (3.915,-1.930) rectangle ++(0.180,0.215);
\path[fill=black!65, draw=black!65] (4.130,-1.930) rectangle ++(0.180,0.215);
\path[fill=black!65, draw=black!65] (4.345,-1.930) rectangle ++(0.180,0.215);
\path[fill=black!65, draw=black!65] (4.560,-1.930) rectangle ++(0.180,0.215);
\path[fill=black!65, draw=black!65] (4.775,-1.930) rectangle ++(0.180,0.215);
\path[fill=black!65, draw=black!65] (4.990,-1.930) rectangle ++(0.180,0.215);
\path[fill=black!65, draw=black!65] (5.205,-1.930) rectangle ++(0.180,0.215);
\path[fill=black!65, draw=black!65] (5.420,-1.930) rectangle ++(0.180,0.215);
\path[fill=black!65, draw=black!65] (5.635,-1.930) rectangle ++(0.180,0.215);
\path[fill=black!65, draw=black!65] (5.850,-1.930) rectangle ++(0.180,0.215);
\path[fill=black!65, draw=black!65] (6.065,-1.930) rectangle ++(0.180,0.215);
\path[fill=black!65, draw=black!65] (6.280,-1.930) rectangle ++(0.180,0.215);
\node[anchor=east] at (1.470,-2.073) {Deopt.};
\path[fill=black!65, draw=black!65] (1.550,-2.180) rectangle ++(0.180,0.215);
\path[fill=deopt, draw=deopt] (1.765,-2.180) rectangle ++(0.180,0.215);
\path[fill=black!65, draw=black!65] (1.980,-2.180) rectangle ++(0.180,0.215);
\path[fill=black!65, draw=black!65] (2.195,-2.180) rectangle ++(0.180,0.215);
\path[fill=black!65, draw=black!65] (2.410,-2.180) rectangle ++(0.180,0.215);
\path[fill=black!65, draw=black!65] (2.625,-2.180) rectangle ++(0.180,0.215);
\path[fill=black!65, draw=black!65] (2.840,-2.180) rectangle ++(0.180,0.215);
\path[fill=black!65, draw=black!65] (3.055,-2.180) rectangle ++(0.180,0.215);
\path[fill=black!65, draw=black!65] (3.270,-2.180) rectangle ++(0.180,0.215);
\path[fill=black!65, draw=black!65] (3.485,-2.180) rectangle ++(0.180,0.215);
\path[fill=black!65, draw=black!65] (3.700,-2.180) rectangle ++(0.180,0.215);
\path[fill=black!65, draw=black!65] (3.915,-2.180) rectangle ++(0.180,0.215);
\path[fill=black!65, draw=black!65] (4.130,-2.180) rectangle ++(0.180,0.215);
\path[fill=black!65, draw=black!65] (4.345,-2.180) rectangle ++(0.180,0.215);
\path[fill=black!65, draw=black!65] (4.560,-2.180) rectangle ++(0.180,0.215);
\path[fill=black!65, draw=black!65] (4.775,-2.180) rectangle ++(0.180,0.215);
\path[fill=black!65, draw=black!65] (4.990,-2.180) rectangle ++(0.180,0.215);
\path[fill=black!65, draw=black!65] (5.205,-2.180) rectangle ++(0.180,0.215);
\path[fill=black!65, draw=black!65] (5.420,-2.180) rectangle ++(0.180,0.215);
\path[fill=black!65, draw=black!65] (5.635,-2.180) rectangle ++(0.180,0.215);
\path[fill=black!65, draw=black!65] (5.850,-2.180) rectangle ++(0.180,0.215);
\path[fill=black!65, draw=black!65] (6.065,-2.180) rectangle ++(0.180,0.215);
\path[fill=black!65, draw=black!65] (6.280,-2.180) rectangle ++(0.180,0.215);
\node[anchor=west, inner sep=0pt, font=\scriptsize\bfseries] at (0,-2.600) {Taxi};
\node[anchor=west, inner sep=0pt, text=black!60] at (1.550,-2.600) {2022};
\node[anchor=west, inner sep=0pt, text=black!60] at (3.915,-2.600) {2023};
\node[anchor=west, inner sep=0pt, text=black!60] at (6.495,-2.600) {2024};
\node[anchor=east] at (1.470,-2.873) {TFDV read};
\path[fill=black!65, draw=black!65] (1.550,-2.980) rectangle ++(0.180,0.215);
\path[fill=black!65, draw=black!65] (1.765,-2.980) rectangle ++(0.180,0.215);
\path[fill=black!65, draw=black!65] (1.980,-2.980) rectangle ++(0.180,0.215);
\path[fill=black!65, draw=black!65] (2.195,-2.980) rectangle ++(0.180,0.215);
\path[fill=black!65, draw=black!65] (2.410,-2.980) rectangle ++(0.180,0.215);
\path[fill=black!65, draw=black!65] (2.625,-2.980) rectangle ++(0.180,0.215);
\path[fill=black!65, draw=black!65] (2.840,-2.980) rectangle ++(0.180,0.215);
\path[fill=black!65, draw=black!65] (3.055,-2.980) rectangle ++(0.180,0.215);
\path[fill=black!65, draw=black!65] (3.270,-2.980) rectangle ++(0.180,0.215);
\path[fill=black!65, draw=black!65] (3.485,-2.980) rectangle ++(0.180,0.215);
\path[fill=black!65, draw=black!65] (3.700,-2.980) rectangle ++(0.180,0.215);
\path[fill=black!65, draw=black!65] (3.915,-2.980) rectangle ++(0.180,0.215);
\path[fill=black!65, draw=black!65] (4.130,-2.980) rectangle ++(0.180,0.215);
\path[fill=black!65, draw=black!65] (4.345,-2.980) rectangle ++(0.180,0.215);
\path[fill=black!65, draw=black!65] (4.560,-2.980) rectangle ++(0.180,0.215);
\path[fill=black!65, draw=black!65] (4.775,-2.980) rectangle ++(0.180,0.215);
\path[pattern=north east lines, pattern color=black, draw=black] (4.990,-2.980) rectangle ++(0.180,0.215);
\path[fill=black!65, draw=black!65] (5.205,-2.980) rectangle ++(0.180,0.215);
\path[pattern=north east lines, pattern color=black, draw=black] (5.420,-2.980) rectangle ++(0.180,0.215);
\path[pattern=north east lines, pattern color=black, draw=black] (5.635,-2.980) rectangle ++(0.180,0.215);
\path[pattern=north east lines, pattern color=black, draw=black] (5.850,-2.980) rectangle ++(0.180,0.215);
\path[pattern=north east lines, pattern color=black, draw=black] (6.065,-2.980) rectangle ++(0.180,0.215);
\path[pattern=north east lines, pattern color=black, draw=black] (6.280,-2.980) rectangle ++(0.180,0.215);
\path[pattern=north east lines, pattern color=black, draw=black] (6.495,-2.980) rectangle ++(0.180,0.215);
\path[fill=black!65, draw=black!65] (6.710,-2.980) rectangle ++(0.180,0.215);
\path[fill=black!65, draw=black!65] (6.925,-2.980) rectangle ++(0.180,0.215);
\path[fill=black!65, draw=black!65] (7.140,-2.980) rectangle ++(0.180,0.215);
\path[fill=black!65, draw=black!65] (7.355,-2.980) rectangle ++(0.180,0.215);
\path[fill=black!65, draw=black!65] (7.570,-2.980) rectangle ++(0.180,0.215);
\path[fill=black!65, draw=black!65] (7.785,-2.980) rectangle ++(0.180,0.215);
\path[fill=black!65, draw=black!65] (8.000,-2.980) rectangle ++(0.180,0.215);
\path[fill=black!65, draw=black!65] (8.215,-2.980) rectangle ++(0.180,0.215);
\path[fill=black!65, draw=black!65] (8.430,-2.980) rectangle ++(0.180,0.215);
\path[fill=black!65, draw=black!65] (8.645,-2.980) rectangle ++(0.180,0.215);
\path[fill=black!65, draw=black!65] (8.860,-2.980) rectangle ++(0.180,0.215);
\node[anchor=east] at (1.470,-3.123) {Full cert.};
\path[fill=black!65, draw=black!65] (1.550,-3.230) rectangle ++(0.180,0.215);
\path[fill=black!65, draw=black!65] (1.765,-3.230) rectangle ++(0.180,0.215);
\path[fill=black!65, draw=black!65] (1.980,-3.230) rectangle ++(0.180,0.215);
\path[fill=black!65, draw=black!65] (2.195,-3.230) rectangle ++(0.180,0.215);
\path[fill=black!65, draw=black!65] (2.410,-3.230) rectangle ++(0.180,0.215);
\path[fill=black!65, draw=black!65] (2.625,-3.230) rectangle ++(0.180,0.215);
\path[fill=black!65, draw=black!65] (2.840,-3.230) rectangle ++(0.180,0.215);
\path[fill=black!65, draw=black!65] (3.055,-3.230) rectangle ++(0.180,0.215);
\path[fill=black!65, draw=black!65] (3.270,-3.230) rectangle ++(0.180,0.215);
\path[fill=black!15, draw=black!15] (3.485,-3.230) rectangle ++(0.180,0.215);
\path[fill=black!65, draw=black!65] (3.700,-3.230) rectangle ++(0.180,0.215);
\path[fill=black!65, draw=black!65] (3.915,-3.230) rectangle ++(0.180,0.215);
\path[fill=black!65, draw=black!65] (4.130,-3.230) rectangle ++(0.180,0.215);
\path[fill=black!65, draw=black!65] (4.345,-3.230) rectangle ++(0.180,0.215);
\path[fill=black!65, draw=black!65] (4.560,-3.230) rectangle ++(0.180,0.215);
\path[fill=black!65, draw=black!65] (4.775,-3.230) rectangle ++(0.180,0.215);
\path[fill=white, draw=black, line width=0.5pt] (4.990,-3.230) rectangle ++(0.180,0.215);
\path[fill=black!65, draw=black!65] (5.205,-3.230) rectangle ++(0.180,0.215);
\path[fill=white, draw=black, line width=0.5pt] (5.420,-3.230) rectangle ++(0.180,0.215);
\path[fill=white, draw=black, line width=0.5pt] (5.635,-3.230) rectangle ++(0.180,0.215);
\path[fill=white, draw=black, line width=0.5pt] (5.850,-3.230) rectangle ++(0.180,0.215);
\path[fill=white, draw=black, line width=0.5pt] (6.065,-3.230) rectangle ++(0.180,0.215);
\path[fill=white, draw=black, line width=0.5pt] (6.280,-3.230) rectangle ++(0.180,0.215);
\path[fill=white, draw=black, line width=0.5pt] (6.495,-3.230) rectangle ++(0.180,0.215);
\path[fill=black!65, draw=black!65] (6.710,-3.230) rectangle ++(0.180,0.215);
\path[fill=black!65, draw=black!65] (6.925,-3.230) rectangle ++(0.180,0.215);
\path[fill=black!65, draw=black!65] (7.140,-3.230) rectangle ++(0.180,0.215);
\path[fill=black!65, draw=black!65] (7.355,-3.230) rectangle ++(0.180,0.215);
\path[fill=black!65, draw=black!65] (7.570,-3.230) rectangle ++(0.180,0.215);
\path[fill=black!65, draw=black!65] (7.785,-3.230) rectangle ++(0.180,0.215);
\path[fill=black!65, draw=black!65] (8.000,-3.230) rectangle ++(0.180,0.215);
\path[fill=black!65, draw=black!65] (8.215,-3.230) rectangle ++(0.180,0.215);
\path[fill=black!65, draw=black!65] (8.430,-3.230) rectangle ++(0.180,0.215);
\path[fill=black!65, draw=black!65] (8.645,-3.230) rectangle ++(0.180,0.215);
\path[fill=black!65, draw=black!65] (8.860,-3.230) rectangle ++(0.180,0.215);
\node[anchor=east] at (1.470,-3.373) {Deopt.};
\path[fill=black!65, draw=black!65] (1.550,-3.480) rectangle ++(0.180,0.215);
\path[fill=black!65, draw=black!65] (1.765,-3.480) rectangle ++(0.180,0.215);
\path[fill=black!65, draw=black!65] (1.980,-3.480) rectangle ++(0.180,0.215);
\path[fill=black!65, draw=black!65] (2.195,-3.480) rectangle ++(0.180,0.215);
\path[fill=black!65, draw=black!65] (2.410,-3.480) rectangle ++(0.180,0.215);
\path[fill=black!65, draw=black!65] (2.625,-3.480) rectangle ++(0.180,0.215);
\path[fill=black!65, draw=black!65] (2.840,-3.480) rectangle ++(0.180,0.215);
\path[fill=black!65, draw=black!65] (3.055,-3.480) rectangle ++(0.180,0.215);
\path[fill=black!65, draw=black!65] (3.270,-3.480) rectangle ++(0.180,0.215);
\path[fill=deopt, draw=deopt] (3.485,-3.480) rectangle ++(0.180,0.215);
\path[fill=black!65, draw=black!65] (3.700,-3.480) rectangle ++(0.180,0.215);
\path[fill=black!65, draw=black!65] (3.915,-3.480) rectangle ++(0.180,0.215);
\path[fill=black!65, draw=black!65] (4.130,-3.480) rectangle ++(0.180,0.215);
\path[fill=black!65, draw=black!65] (4.345,-3.480) rectangle ++(0.180,0.215);
\path[fill=black!65, draw=black!65] (4.560,-3.480) rectangle ++(0.180,0.215);
\path[fill=black!65, draw=black!65] (4.775,-3.480) rectangle ++(0.180,0.215);
\path[fill=white, draw=black, line width=0.5pt] (4.990,-3.480) rectangle ++(0.180,0.215);
\path[fill=black!65, draw=black!65] (5.205,-3.480) rectangle ++(0.180,0.215);
\path[fill=white, draw=black, line width=0.5pt] (5.420,-3.480) rectangle ++(0.180,0.215);
\path[fill=white, draw=black, line width=0.5pt] (5.635,-3.480) rectangle ++(0.180,0.215);
\path[fill=white, draw=black, line width=0.5pt] (5.850,-3.480) rectangle ++(0.180,0.215);
\path[fill=white, draw=black, line width=0.5pt] (6.065,-3.480) rectangle ++(0.180,0.215);
\path[fill=white, draw=black, line width=0.5pt] (6.280,-3.480) rectangle ++(0.180,0.215);
\path[fill=white, draw=black, line width=0.5pt] (6.495,-3.480) rectangle ++(0.180,0.215);
\path[fill=black!65, draw=black!65] (6.710,-3.480) rectangle ++(0.180,0.215);
\path[fill=black!65, draw=black!65] (6.925,-3.480) rectangle ++(0.180,0.215);
\path[fill=black!65, draw=black!65] (7.140,-3.480) rectangle ++(0.180,0.215);
\path[fill=black!65, draw=black!65] (7.355,-3.480) rectangle ++(0.180,0.215);
\path[fill=black!65, draw=black!65] (7.570,-3.480) rectangle ++(0.180,0.215);
\path[fill=black!65, draw=black!65] (7.785,-3.480) rectangle ++(0.180,0.215);
\path[fill=black!65, draw=black!65] (8.000,-3.480) rectangle ++(0.180,0.215);
\path[fill=black!65, draw=black!65] (8.215,-3.480) rectangle ++(0.180,0.215);
\path[fill=black!65, draw=black!65] (8.430,-3.480) rectangle ++(0.180,0.215);
\path[fill=black!65, draw=black!65] (8.645,-3.480) rectangle ++(0.180,0.215);
\path[fill=black!65, draw=black!65] (8.860,-3.480) rectangle ++(0.180,0.215);
\matrix[anchor=north west, column sep=2pt, ampersand replacement=\&] at (1.550,-3.650) {
\path[fill=black!65, draw=black!65] (0,0) rectangle ++(0.180,0.215); \& \node[anchor=west, inner sep=1pt] {ran, wrote the fresh compilation's files}; \& \path[fill=deopt, draw=deopt] (0,0) rectangle ++(0.180,0.215); \& \node[anchor=west, inner sep=1pt] {ran with removals re-enabled, as fresh}; \& \path[fill=white, draw=black, line width=0.5pt] (0,0) rectangle ++(0.180,0.215); \& \node[anchor=west, inner sep=1pt] {stopped where the fresh compilation stops}; \& \path[fill=black!15, draw=black!15] (0,0) rectangle ++(0.180,0.215); \& \node[anchor=west, inner sep=1pt] {refused although fresh runs}; \& \path[pattern=north east lines, pattern color=black, draw=black] (0,0) rectangle ++(0.180,0.215); \& \node[anchor=west, inner sep=1pt] {wrong output}; \\
};
\end{tikzpicture}
\end{adjustbox}
\caption{RQ4, compiled on the first month: every later month on a calendar axis (the traffic data has a gap), for the best validator, the full certificate and selective deoptimization. White: a fresh compilation stops too (temperatures of 0\,K in 2014, taxi fares above \$1,000). The grey run in the taxi rows begins in February 2023, when the publisher changed two column types; TFDV admits those months, and the optimized program writes them with the old types.}
\label{fig:timeline}
\end{figure*}


Over \RealRuns{} runs (Table~\ref{tab:rq-real}, Figs.~\ref{fig:policies} and~\ref{fig:timeline}), selective deoptimization
handles \RealDeoptCorrect{} (\RealDeoptCorrectPct\%) exactly as a fresh compilation would, with
\RealDeoptUnsound{} wrong; the full certificate handles \RealFullCorrect{}. The difference is
\RealRescued{} months the full certificate refuses. In three, the traffic sensor reports temperatures
of 0\,K; the re-enabled range check stops the run, as the fresh compilation fails. In four, a month has
an hour with 9,831\,mm of rain or a day with humidity 0, which the compiled month had not; the removed
cleaning filter runs again and the run writes what the fresh compilation writes. In two, a taxi trip
has a fare of \$2,020; the re-enabled check runs after the filter that removes that trip, passes, and
the run finishes. All \RealDeoptNeedless{} remaining refusals are one event: in February 2023 the
taxi publisher changed \texttt{passenger\_count} and \texttt{payment\_type} from double to integer,
and the program compiled on January 2022 reads them as \textsc{Float}. A type is a fact no run-time
code of the compiled program stands in for; monthly compilation follows the change after one month.

The validators are either strict or unsound. Great Expectations' profile constrains means, medians,
quantiles and unique counts, and pandera's inferred schema the compiled month's range of every
column; set up on every column or only on the columns read, they refuse
\RealValidatorStopMin--\RealValidatorStopMax\% of the months although a fresh compilation runs them.
TFDV infers types, presence and string domains but no numeric ranges. On every column it refuses
\RealTfdvAllNeedlessPct\% of the months, for string values of columns the program never reads (a new
holiday name refuses a traffic month); on the columns read it refuses fewer, but both setups admit
months on which the optimized program does not do what a fresh compilation does
(\RealTfdvReadUnsound{} months): taxi months after the type change, written with the old column types,
and months on which a removed range check should have stopped the run, taxi fares above \$1,000 and
traffic temperatures of 0\,K, where the optimized program finishes and writes the bad rows. Generic
validators do not know which facts the compiled program relies on.

\subsection{RQ5: Real pandas scripts}\label{sec:rq-corpus}

\paragraph{Corpora.} GitHub code search for Python files containing \texttt{read\_csv},
\texttt{fillna} and \texttt{Embarked} (the Kaggle Titanic data) returned \CorpusAll{} distinct files,
and for \texttt{read\_csv}, \texttt{fillna} and \texttt{LotFrontage} (Kaggle House Prices) \HousesAll{}
files. We used the first 100 Titanic scripts as a development set, to choose which idioms the lifter
handles; the remaining \CorpusTest{} and all House Prices scripts are held out: we did not look at them
before the final run, and the lifter was not changed for House Prices. A script is in scope if it reads
the training file with \texttt{pd.read\_csv(path)} at module level. The House Prices files mark missing
values with the text \texttt{NA}, and one column has a category called \texttt{None}; pandas reads both
as missing. We give both sides the same files, with an empty cell, \PL's notation, wherever pandas reads
a missing value. Our first run replaced only \texttt{NA}: two scripts then differed from pandas on every
batch, the training file included, because \PL read \texttt{None} as a value that the scripts fill.
That is a difference of CSV notation, independent of reuse; an exact lifter must carry pandas' list of
missing-value strings into \texttt{LOAD}.

\paragraph{Protocol.} The lifted program is compiled against the training file and given eleven new
batches: Kaggle's real test file (no target column, and missing values in columns where the training
file has none), five bootstrap resamples of the training rows, and five copies with 20\% extra missing
values in one column. Each batch is run by the compiled program and, as ground truth, by the lifted
pandas statements; the cleaned tables are compared value by value.

\begin{table}[t]
\caption{Real pandas scripts from GitHub (RQ5): held-out scripts, lifted, compiled once on the training file and run on new batches; the ground truth is the lifted pandas statements run on the same batch.}
\label{tab:rq-corpus}
\centering
\footnotesize
\begin{adjustbox}{max width=\linewidth}
\begin{tabular}{@{}lrr@{}}
\toprule
& Titanic & House Prices \\
\midrule
Held-out scripts & 272 & 461 \\
Lifted, in scope & 130 & 168 \\
Compiled & 124 & 168 \\
Runs on new batches & 1364 & 1848 \\
\emph{Full certificate:} admitted & 1347 & 1838 \\
\quad identical to the script & 1347 & 1838 \\
\quad refused, script runs & 10 & 3 \\
\emph{With deoptimization:} admitted & 1357 & 1841 \\
\quad identical to the script & 1357 & 1841 \\
\quad stopped by a check & 0 & 0 \\
\quad different (unsound) & 0 & 0 \\
\quad refused, script runs & 0 & 0 \\
\quad ran with removals re-enabled & 10 & 3 \\
\bottomrule
\end{tabular}
\end{adjustbox}
\end{table}


\paragraph{Results.} With deoptimization, \CorpusSame{} Titanic runs and \HousesSame{} House Prices
runs reproduce the scripts' results exactly, and none differ (Table~\ref{tab:rq-corpus}). The
difference to the full certificate is the price of optimizing for the compiled data: the training
files have no missing \texttt{Fare} (Titanic) or \texttt{MSZoning} and \texttt{GarageCars} (House
Prices), so the compiler removes the scripts' \texttt{fillna} on them, and the full certificate then
refuses every batch with missing values there, including Kaggle's own test files. Deoptimization runs
the removed \texttt{fillna} again for exactly those batches (\CorpusDeopted{} Titanic and
\HousesDeopted{} House Prices runs). Most House Prices scripts fill columns that already have missing
values in the training file, such as \texttt{LotFrontage}; those fills are never removed, and the
certificate does not constrain them. The remaining refusals are batches the scripts crash on too:
Kaggle's test files have no target column. The second lifter version (Section~\ref{sec:lift}), written
on other scripts after these corpora were evaluated, lifts \CorpusBothDerStatements{} cleaning
statements of the same held-out scripts instead of \CorpusBothStatements{}; the runs end as before
(\CorpusBothDerSame{} of \CorpusBothDerRuns{} identical to the scripts, \CorpusBothDerUnsound{}
different).

\begin{table}[t]
\caption{Other people's pandas scripts for the monthly traffic and bike data (RQ5): lifted with derived columns and dates, compiled on the first month and run on every later month; the ground truth is the lifted pandas statements run on the same month. The traffic scripts were the development set of the second lifter version.}
\label{tab:rq-drift}
\centering
\footnotesize
\setlength{\tabcolsep}{4pt}
\begin{adjustbox}{max width=\linewidth}
\begin{tabular}{@{}lrrr@{}}
\toprule
 & Traffic & Bikes, daily & Bikes, hourly \\
 & (development) & (held out) & (held out) \\
\midrule
GitHub files & 241 & 254 & 450 \\
lifted, in scope & 47 & 26 & 14 \\
compiled on the first month & 45 & 26 & 14 \\
\quad with a removal that can run again & 3 & 0 & 1 \\
statements lifted & 110 & 49 & 23 \\
\quad of them \texttt{DERIVE} & 52 & 26 & 7 \\
\quad of them \texttt{PARSE\_DATE} & 41 & 15 & 8 \\
runs on later months & 2,790 & 598 & 322 \\
\emph{With deoptimization:} as the script & 2,790 & 598 & 322 \\
\quad stopped where the script fails & 0 & 0 & 0 \\
\quad refused, script runs & 0 & 0 & 0 \\
\quad different (wrong) & 0 & 0 & 0 \\
\quad ran with removals re-enabled & 6 & 0 & 0 \\
\emph{Full certificate:} as the script & 2,784 & 598 & 322 \\
\quad refused, script runs & 6 & 0 & 0 \\
\quad different (wrong) & 0 & 0 & 0 \\
\bottomrule
\end{tabular}
\end{adjustbox}
\end{table}


\paragraph{Other people's pipelines on real monthly data.} The pipelines of RQ4 are ours, and most
Kaggle batches above are made, not observed. GitHub code search for Python files that read the UCI
traffic file, or the bike-sharing \texttt{day.csv} or \texttt{hour.csv}, returned \DriftAll{} files.
Their cleaning mostly begins by parsing the date column and deriving hours, weekdays or Celsius
temperatures from it, which the first lifter version cannot express. With \texttt{PARSE\_DATE} and
\texttt{DERIVE} in the language and the second lifter version, developed on the traffic scripts while
the bike scripts stayed held out, \DriftCompiled{} scripts compile on the first month of the real data
(\DriftDerive{} \texttt{DERIVE} and \DriftParseDate{} \texttt{PARSE\_DATE} among \DriftStatements{}
lifted statements) and run on every later month (Table~\ref{tab:rq-drift}). All \DriftRuns{} runs end
as the scripts' own statements do on the same month, and none is refused. Most of these programs have
no proof that depends on values that drift: their certificates require column types, and for
\DriftDatesAdmitted{} of them that a column holds dates; \DriftWithRemovals{} have a removal that can run
again. Two traffic scripts filter impossible readings themselves (\texttt{temp > 0},
\texttt{rain\_1h < 500}). On October 2012 both filters keep every row and are removed, and they run
again in exactly the three months with a temperature of 0\,K or 9,831\,mm of rain, where the full
certificate refuses the month (\DriftTrafficFullNeedless{} runs). One lifted script fills a text column
with the number 0, which pandas accepts and \PL rejects as a type error.

\subsection{RQ6: Composing stages}\label{sec:rq-compose}

Each workload of RQ4 gets a second stage that reads the cleaned table, as a feature step would: range
and not-null contracts, a fill, a filter and a normalization. Both stages are compiled once (stage~2 on
what stage~1 writes for the anchor month), and \texttt{tools/contracts.py} compares stage~1's guarantee
with what stage~2 relies on. Every later month runs through both compiled stages, and stage~2 is also
compiled and run fresh on stage~1's output.

\begin{table}[t]
\caption{Composing two compiled stages on real monthly data (RQ6). Facts: what stage~2 relies on (admission and removals); Guar.: implied by stage~1's guarantee; Months: run through both stages; Deopt.: months stage~2 deoptimized; by a guaranteed fact: never.}
\label{tab:rq-compose}
\centering
\footnotesize
\begin{adjustbox}{max width=\linewidth}
\begin{tabular}{@{}llrrrrr@{}}
\toprule
Data & Stage 1 & Facts & Guar. & Months & Deopt. & Guar.\ broken \\
\midrule
taxi & as written & 13 & 11 & 15 & 0 & 0 \\
taxi & strengthened & 13 & 13 & 15 & 0 & 0 \\
traffic & as written & 14 & 12 & 60 & 0 & 0 \\
traffic & strengthened & 14 & 14 & 60 & 0 & 0 \\
bikes & as written & 11 & 10 & 23 & 0 & 0 \\
bikes & strengthened & 11 & 11 & 23 & 0 & 0 \\
\bottomrule
\end{tabular}
\end{adjustbox}
\end{table}


The composition shows that stage~1 already guarantees most of what stage~2 relies on
(\ComposeGuaranteed{} of \ComposeFacts{} facts across the three workloads), so stage~2's checks of
those facts can never fire on stage~1's output (Table~\ref{tab:rq-compose}). The facts it reports as
open are genuine gaps between the stages: stage~1 never bounds taxi tips, which stage~2 expects in
$[0,500]$; it fills missing rain with 0 but never bounds rain from below; and it never checks the
traffic cloud cover or the bike temperatures for missing values, which stage~2's filters, removed as
always true, rely on. For each, the tool wrote the \texttt{EXPECT} that stage~1 would need. Across
\ComposeMonths{} months through both stages (in both versions), no fact reported as guaranteed ever
made stage~2 refuse or deoptimize. The open facts also held in these months, so stage~2 never
deoptimized at all; the composition's value here is that after the suggested statements are added to
stage~1, every fact stage~2 relies on is guaranteed statically, for all future months, and adding
them made no month stop earlier in stage~1 (the same \ComposeStageOneStops{} months stop there in both
versions, all at the checks stage~1 already had). Every stage-2 run ended exactly like a fresh
compilation of stage~2 on the same input.

\subsection{RQ7: The cost of certification and reuse}\label{sec:rq-cost}

Reusing a compiled program costs its admission check and, with deoptimization, the facts of the columns
item certificates mention. Over the \OhMonths{} real months on which every variant finished (median of
three runs each), the median run took \OhTrusted\,s with no check at all, \OhFull\,s with admission
against the full certificate, and \OhDeopt\,s with admission and deoptimization. Admission with
deoptimization changed the run time by \OhPct\% at the median (interquartile range \OhPctLow\% to
\OhPctHigh\%), which is inside the run-to-run noise of the machine: the check is one pass over a few
columns the program reads anyway. Compiling a month afresh took \OhCompileMs\,ms before running and \OhFresh\,s with the run.

Certification is paid once, at compile time. Over \CcPrograms{} programs (the shipped examples and
the three real pipelines), the compiler ran \CcAnalysesMin{} to \CcAnalysesMax{} re-analyses with
deoptimization against at most \CcFullAnalysesMax{} for the full certificate alone (at most
\CcRatioMax$\times$ as many), and every compilation took \CcMsMin{} to \CcMsMax\,ms. The analysis is
cheap because it runs on dataset states, not on data.

The certificates are small enough to read. For the \CsPrograms{} table programs, counting a type, the
absence of missing values and each bound as one fact, the full certificate has \CsFullMin{} to
\CsFullMax{} facts, the admission certificate here holds only the type of each certified column, and
each of the \CsItems{} item certificates needs \CsItemMin{} or \CsItemMax{} facts beyond it. The real
pipelines' certificates mention \CsTrafficCols{} of the \CsTrafficFileCols{} traffic columns,
\CsBikesCols{} of \CsBikesFileCols{} bike columns and \CsTaxiCols{} of \CsTaxiFileCols{} taxi columns,
with bounds only at the program's constants. The validators constrain every column they are set up on
with facts no proof needs, such as the compiled month's exact range or mean; that, not the number of
columns, is why Great Expectations and pandera refuse as often on the columns read as on all of them.

We state the consequence plainly: at these sizes reuse does not save time, because a fresh compilation is
itself cheap. What reuse buys is a fixed, reviewed artifact that runs unchanged on every admitted batch
(\S\ref{sec:rq-real}), an exact signal when a batch drifts, naming the removal it breaks, and output
guarantees a downstream stage can rely on without seeing the data (\S\ref{sec:rq-compose}). Recompiling
gives none of these: each month gets a different program and no statement about what it assumed.

\subsection{RQ8: Compiling against data you cannot see}\label{sec:rq-private}

\paragraph{Protocol.} For each RQ4 pipeline (compiled for its first month) and each of the
\PrivTitanicScripts{} RQ5 Titanic scripts that compile (compiled for the training file), the data owner
ran \texttt{--queries} and released facts about the compiled file, and the program was compiled with
\texttt{--facts} in a folder that holds only the program and the facts file. The facts came from the
release at the program's cells with $\varepsilon \in \{0.1, 1, 10, 100\}$ ($\varepsilon$ is the budget of
each table a program loads, $\beta = 0.05$; five seeds for the pipelines, one for the scripts, so the
script rows spend one release per script on the same training file); from one release
for all Titanic scripts together (two seeds); from the same cells without noise; and from four baselines
that do not use the program's numbers: \emph{generic} bins (the same budget per column on a fixed
logarithmic grid, $[0.001, 0.01)$, \ldots, $[10^5, 10^6)$, the negative ones and the outer ones, read the
same way), \emph{synthetic} data (a table sampled from the generic release, values uniform in occupied
bins, compiled normally, as an analyst would who develops against differentially private synthetic or
dummy data), \emph{optimistic} facts (no value is missing, chosen without the data), and no facts at all.
Random facts, true or false at random, test exactness. Every generated program was run on every batch of
RQ4 or RQ5 and on the compiled file itself, and scored against a fresh compilation of the batch.

\paragraph{Exactness.} \PrivCompilations{} compilations ran \PrivRuns{} times; \PrivUnsound{} runs were
unsound or stopped wrongly, for every kind of facts, including the \PrivRandom{} compilations from random
facts (\PrivRandomRuns{} runs). All \PrivNeedless{} refusals of batches a fresh compilation runs are taxi
months with the schema change of February 2023, which the normal compilation refuses too
(\S\ref{sec:rq-real}). The unit test of the private compilation checks the same property on
\PrivTestFacts{} random facts files with random inputs, \PrivTestDeopt{} of which deoptimize.

\paragraph{What the analyst learns.} Table~\ref{tab:rq-private} compares the removals of each compilation
with those of the normal compilation of the same file. Without noise, the program's cells give exactly the
normal removals (recall \PrivExactRecall\%): the analysis compares bounds only with those numbers. Recall
alone is cheap: the optimistic facts reach \PrivOptimisticRecall\% of it without data, at a precision of
\PrivOptimisticPrecision\%. With the program's cells at $\varepsilon = 1$, the taxi and traffic pipelines
keep all normal removals at full precision, while the generic bins, at the same budget, keep
\PrivGenTaxiEpsOneRecall\% (taxi) and \PrivGenTrafficEpsOneRecall\% (traffic): the program's numbers
(\texttt{BETWEEN 1 AND 265}, \texttt{temp BETWEEN 230 AND 320}) do not fall on the grid. Synthetic data
compiled normally is worse: none of the \PrivSynEpsOneFailed{} compilations of the three pipelines from
synthetic data succeeds, at $\varepsilon = 1$ or $10$; each stops with contracts the real data satisfies
(a temperature of 867~K, sampled uniformly in its occupied bin), so the analyst gets errors where the
private compilation gives a program. What can be
learned depends on how many rows a fact concerns: a median month has \PrivTrafficRows{} traffic rows,
\PrivBikesRows{} bike rows and \PrivTaxiRows{} taxi trips, and the bike months, one row a day, keep the
range removals only at $\varepsilon = 10$. On the Titanic training file, removal precision rises from
\PrivProgTitanicEpsOnePrecision\% at $\varepsilon = 1$ to \PrivProgTitanicEpsTenPrecision\% at
$\varepsilon = 10$ (generic bins and synthetic data: \PrivGenTitanicEpsTenPrecision\%); the false facts
are mostly ``no missing \texttt{Embarked}'', whose two missing values no release can reveal, and the run on
the training file deoptimizes them. One release for all \PrivTitanicScripts{} scripts at a total of
$\varepsilon = 1$ gives each of them \PrivCorpTitanicEpsOnePrecision\% precision, about what a script gets
from its own release at $\varepsilon = 1$ (\PrivProgTitanicEpsOnePrecision\%), which together spend
\PrivTitanicScripts{} times the budget: the programs share the cells of the columns they all use.
Protecting tickets instead of rows (\PrivTicketGroups{} of the \PrivTicketCount{} tickets of the training file
have several passengers, up to \PrivTicketMax{}), with one row counted per ticket, gives
\PrivUnitTitanicEpsOnePrecision\% precision at $\varepsilon = 1$, as rows do, and
\PrivUnitTitanicEpsTenPrecision\% at $\varepsilon = 10$, where rows give \PrivProgTitanicEpsTenPrecision\%: the
tables can no longer share the budget in parallel, and the rows not counted make some facts false, which
the runs deoptimize. A family on one ticket is then protected at $\varepsilon$ rather than at up to
\PrivTicketMax$\varepsilon$, and every run stays exact.

\paragraph{Warnings.} On a \PrivDiagRows-row Parquet file we injected a mistake into $m$ rows (amounts
above the contract's bound, missing values in a \texttt{NOT\_NULL} column, values below a bound)
(Table~\ref{tab:rq-private-diag}). A normal compilation reports the broken contract for every $m > 0$. A
private compilation from the program's cells warns that it is likely violated in every run from
$m = \PrivDiagEpsOne$ at $\varepsilon = 1$ and $m = \PrivDiagEpsTen$ at $\varepsilon = 10$, with
\PrivDiagFalse{} warnings about a contract the data does not break in \PrivDiagCompilations{} compilations,
as Lemma~\ref{lem:utility} predicts. Generic bins, and synthetic data run on the program (as on a
validation server, where a run stops at its first failing check), report every mistake in every run from
the same sizes ($m = \PrivDiagGenericEpsOne$ at $\varepsilon = 1$), also without false reports
(\PrivDiagGenericFalse{} and \PrivDiagSyntheticFalse{}); between those sizes they report less often: at
$\varepsilon = 0.3$, a mistake in 100 rows in 50--60\% of the runs where the program's cells report it in
80--100\%, and at $\varepsilon = 3$, one in 10 rows in 40--60\% against 80--90\%. They spend the budget on
about three times as many cells; the injected values fall in the grid's decades, so the alignment that
decides the removals above does not matter here. A mistake in one row is never revealed by any private
method; the run-time check stops the program at it.

\begin{table*}[t]
\caption{Private compilation (RQ8). Per workload, recall/precision (\%) of the compilation's removals against the normal compilation of the same file: recall is the share of normal removals it also makes, precision the share of its removals that the normal compilation also makes. Compiled: compilations without errors. Deopt.\ on file: compilations whose run on the compiled file itself deoptimizes (a released fact was false). Every run is scored against a fresh compilation of its batch. $\varepsilon$ is the budget of each table of one program; ``one release'' serves all Titanic scripts with one budget. HDMM: the high-dimensional matrix mechanism given the program's comparisons as its workload; AIM: at $(\varepsilon, 10^{-9})$ over the program's cells, read with the release's threshold from its model's marginals, read off its synthetic rows, or its synthetic table compiled normally. Generic: fixed logarithmic bins; synthetic: a table sampled from them, compiled normally.}
\label{tab:rq-private}
\centering
\footnotesize
\begin{adjustbox}{max width=\linewidth}
\begin{tabular}{@{}lrrrrrrrrr@{}}
\toprule
Facts & Traffic & Bikes & Taxi & Titanic & Compilations & Compiled & Deopt.\ on file & Runs & Wrong \\
\midrule
program's cells, no noise & 100/100 & 100/100 & 100/100 & 100/100 & 187 & 187 & 0\% & 2,331 & 0 \\
program's cells, $\varepsilon = 0.1$ & 100/100 & 14/100 & 100/100 &  & 15 & 15 & 0\% & 615 & 0 \\
program's cells, $\varepsilon = 1$ & 100/100 & 14/100 & 100/100 & 100/17 & 199 & 199 & 37\% & 2,823 & 0 \\
program's cells, $\varepsilon = 10$ & 100/100 & 100/100 & 100/100 & 100/31 & 199 & 199 & 16\% & 2,823 & 0 \\
program's cells, $\varepsilon = 100$ & 100/100 & 100/100 & 100/100 &  & 15 & 15 & 0\% & 615 & 0 \\
one release, $\varepsilon = 1$ &  &  &  & 100/17 & 368 & 368 & 40\% & 4,416 & 0 \\
one release, $\varepsilon = 10$ &  &  &  & 100/17 & 368 & 368 & 40\% & 4,416 & 0 \\
per ticket, $\varepsilon = 1$ &  &  &  & 100/17 & 184 & 184 & 40\% & 2,208 & 0 \\
per ticket, $\varepsilon = 10$ &  &  &  & 100/17 & 184 & 184 & 40\% & 2,208 & 0 \\
HDMM, $\varepsilon = 1$ & 100/100 & 14/100 & 100/100 & 100/17 & 199 & 199 & 37\% & 2,823 & 0 \\
HDMM, $\varepsilon = 10$ & 100/100 & 100/100 & 100/100 & 100/20 & 199 & 199 & 29\% & 2,823 & 0 \\
AIM, marginals, $\varepsilon = 1$ & 100/100 & 14/100 & 100/100 & 100/16 & 199 & 199 & 37\% & 2,823 & 0 \\
AIM, marginals, $\varepsilon = 10$ & 100/100 & 77/100 & 100/100 & 100/18 & 199 & 199 & 33\% & 2,823 & 0 \\
AIM, synthetic rows, $\varepsilon = 1$ & 40/100 & 43/100 & 52/100 & 38/21 & 199 & 199 & 12\% & 2,823 & 0 \\
AIM, synthetic rows, $\varepsilon = 10$ & 66/100 & 66/100 & 12/100 & 31/38 & 199 & 199 & 4\% & 2,823 & 0 \\
AIM, compiled, $\varepsilon = 1$ & --/-- & --/-- & --/-- & 44/23 & 199 & 184 & 13\% & 2,208 & 0 \\
AIM, compiled, $\varepsilon = 10$ & --/-- & 100/100 & --/-- & 31/38 & 199 & 185 & 4\% & 2,232 & 0 \\
generic bins, $\varepsilon = 1$ & 83/100 & 14/100 & 60/100 & 100/17 & 199 & 199 & 37\% & 2,823 & 0 \\
generic bins, $\varepsilon = 10$ & 71/100 & 86/100 & 60/100 & 100/20 & 199 & 199 & 30\% & 2,823 & 0 \\
synthetic data, $\varepsilon = 1$ & --/-- & --/-- & --/-- & 100/17 & 199 & 184 & 40\% & 2,208 & 0 \\
synthetic data, $\varepsilon = 10$ & --/-- & --/-- & --/-- & 100/20 & 199 & 184 & 32\% & 2,208 & 0 \\
optimistic (no data) & 29/100 & 14/100 & 40/50 & 100/9 & 187 & 187 & 46\% & 2,331 & 0 \\
random (true or false) & 26/100 & 30/100 & 56/64 & 44/21 & 398 & 398 & 13\% & 5,646 & 0 \\
none (claims nothing) & 0/-- & 0/-- & 0/-- & 0/-- & 187 & 187 & 0\% & 2,331 & 0 \\
\bottomrule
\end{tabular}
\end{adjustbox}
\end{table*}

\begin{table}[t]
\caption{Private compilation (RQ8): share of private compilations that warn that the contract a mistake in $m$ of 1{,}000{,}000 rows breaks is likely violated (five seeds; one without noise).}
\label{tab:rq-private-diag}
\centering
\footnotesize
\begin{tabular}{@{}lrrrrr@{}}
\toprule
Mistake & $m$ & $\varepsilon = 0.1$ & $\varepsilon = 1$ & $\varepsilon = 10$ & no noise \\
\midrule
amount above 999 & 0 & 0\% & 0\% & 0\% & 0\% \\
amount above 999 & 1 & 0\% & 0\% & 0\% & 100\% \\
amount above 999 & 10 & 0\% & 0\% & 0\% & 100\% \\
amount above 999 & 100 & 0\% & 0\% & 100\% & 100\% \\
amount above 999 & 1,000 & 0\% & 100\% & 100\% & 100\% \\
amount above 999 & 10,000 & 100\% & 100\% & 100\% & 100\% \\
s missing & 0 & 0\% & 0\% & 0\% & 0\% \\
s missing & 1 & 0\% & 0\% & 0\% & 100\% \\
s missing & 10 & 0\% & 0\% & 0\% & 100\% \\
s missing & 100 & 0\% & 0\% & 100\% & 100\% \\
s missing & 1,000 & 0\% & 100\% & 100\% & 100\% \\
s missing & 10,000 & 100\% & 100\% & 100\% & 100\% \\
a below 0 & 0 & 0\% & 0\% & 0\% & 0\% \\
a below 0 & 1 & 0\% & 0\% & 0\% & 100\% \\
a below 0 & 10 & 0\% & 0\% & 0\% & 100\% \\
a below 0 & 100 & 0\% & 0\% & 100\% & 100\% \\
a below 0 & 1,000 & 0\% & 100\% & 100\% & 100\% \\
a below 0 & 10,000 & 100\% & 100\% & 100\% & 100\% \\
\bottomrule
\end{tabular}
\end{table}


\subsection{Testing soundness}\label{sec:soundtest}

The theorems rest on monotonicity and soundness of the analysis, which are mechanized for the core
of the language and tested for the implementation. Randomized property tests generate tables (with missing, out-of-range, pruned and
dropped columns) and image folders, decide admission and deoptimization with the certificates, and
compile each input afresh: every input the full certificate admits keeps every proof, and every input
admitted with deoptimization compiles without new errors and keeps every removal it does not
deoptimize (or stops at a re-enabled check exactly where the fresh compilation reports a violation).
For the running example, \RandFull{} of 400 random tables fit the full certificate and \RandDeopt{} are
admitted with deoptimization; \RandFolders{} random image folders are admitted with deoptimization,
all keeping their removals. A program with \texttt{PARSE\_DATE}, \texttt{DERIVE} and a filter on the
month admits \RandDates{} of 400 random batches (dates across three years, text that is not a date,
divisors of 0), \RandDatesDeopted{} of them deoptimizing something, none losing a removal it keeps.

\paragraph{What the evaluation found in our implementation.} Each of the following was a case where
the implementation did not follow the definitions; each is fixed and has a regression test. (1)~The
run-time \texttt{MATCH\_RATE} check, always eliminated before certificates existed, ran for the first
time on changed labels and counted broken images that \texttt{REMOVE\_CORRUPT} had removed. (2)~The
first certificate covered only the columns the optimized program loads, so a proven \texttt{NOT\_NULL}
on a column pruned afterwards was unprotected. (3)~Admission let \textsc{Integer} data into a
\textsc{Float} column and an empty column into any type; the property test found inputs whose fresh
compilation decided differently. (4)~On the held-out Titanic set, scripts drop \texttt{Survived} before
any use; with that \texttt{DROP} removed, admission treated the column as optional and admitted
Kaggle's test file, on which the script fails. (5)~The lifter translated a \texttt{drop} that pandas
rejects. (6)~Private compilations from false facts (RQ8) made a lifted Titanic script refuse every
batch: a fill of \texttt{Embarked}, dropped right after, was removed both as a dead store and as having
no effect, and the certifier kept ``no missing \texttt{Embarked}'' as an admission fact of a column the
program never loads. A dead store is removed whatever the data, so its no-effect fact is no longer an
obligation; a normal compilation of a file without missing \texttt{Embarked} had the same needless
refusals waiting for the next batch. The earlier RQs were measured before this fix, which only removes
admission facts.

\subsection{Threats to validity}

The error messages of RQ1 were chosen by the author and one person classified the issues; the
mutants and the script-style interpreter are ours. Performance numbers come from one machine. The
pipelines of RQ4 and RQ6 are ours, written once before running them; the months are not, and RQ5 runs
other people's pipelines on the same months. Validators are
used with their automatic profilers, as a new user would; a hand-tuned suite can do better, but then
encodes what the certificate derives. RQ5 compares with the lifted statements, not the whole scripts,
and the lifted prefixes are short. Metadata-only soundness assumes truthful Parquet statistics. The
mechanization covers a core of the language; the C++ implementation is tested, not proven. The language has
no joins, loops or user functions. Certificates are facts about input columns: a fact needed only
after a filter is still required of every input row. Private compilation treats column names and types
as public; its guarantee covers what the compiler writes, not what the data owner's runs print or how
long they take (a deoptimization message names a row), which stay with the owner.
